\textbf{Phần 1. Trắc nghiệm lựa chọn}
\setcounterpageref{ex}{0}

\begin{ex} %[4H6V27]
Hai đường thẳng vuông góc với nhau tạo thành mấy góc vuông?
\choice
{$1$ góc vuông}
{\True $4$ góc vuông}
{$2$ góc vuông}
{$3$ góc vuông}
\loigiai{
Hai đường thẳng vuông góc cắt nhau tạo thành 4 góc vuông.\\
\textbf{Đáp án của câu hỏi là B}
}
\end{ex}

\begin{ex} %[4H6V27]
Dụng cụ thường dùng để kiểm tra hai đường thẳng có vuông góc với nhau hay không là
\choice
{\True Ê-ke}
{Thước dây}
{Com-pa}
{Thước cuộn}
\loigiai{
Ê-ke có 1 góc vuông dùng để kiểm tra hai đường thẳng vuông góc.\\
\textbf{Đáp án của câu hỏi là A}
}
\end{ex}

\begin{ex} %[4H6V27]
Trong hình chữ nhật $ABCD$, cặp cạnh nào sau đây vuông góc với nhau?
\choice
{$AB$ và $CD$}
{$AD$ và $BC$}
{\True $AB$ và $BC$}
{$AC$ và $BD$}
\loigiai{
Trong hình chữ nhật, hai cạnh kề nhau $AB$ và $BC$ vuông góc với nhau.\\
\textbf{Đáp án của câu hỏi là C}
}
\end{ex}

\begin{ex} %[4H6V27]
Trong hình vuông $MNPQ$, có bao nhiêu cặp cạnh kề nhau vuông góc với nhau?
\choice
{$2$ cặp}
{$3$ cặp}
{$1$ cặp}
{\True $4$ cặp}
\loigiai{
Có 4 cặp cạnh vuông góc tại 4 đỉnh $M, N, P, Q$.\\
\textbf{Đáp án của câu hỏi là D}
}
\end{ex}

\begin{ex} %[4H6V27]
Hai đường chéo của hình thoi có tính chất nào sau đây?
\choice
{Song song với nhau}
{\True Vuông góc với nhau}
{Không cắt nhau}
{Bằng nhau}
\loigiai{
Hai đường chéo của hình thoi vuông góc với nhau.\\
\textbf{Đáp án của câu hỏi là B}
}
\end{ex}

\begin{ex} %[4H6V27]
Cho tam giác $ABC$ có góc $A$ là góc vuông. Cặp cạnh vuông góc với nhau là
\choice
{\True $AB$ và $AC$}
{$AB$ và $BC$}
{$AC$ và $BC$}
{Không có}
\loigiai{
Hai cạnh góc vuông là $AB$ và $AC$.\\
\textbf{Đáp án của câu hỏi là A}
}
\end{ex}

\begin{ex} %[4H6V27]
Hình chữ nhật $ABCD$ có mấy góc vuông?
\choice
{$2$ góc vuông}
{$3$ góc vuông}
{\True $4$ góc vuông}
{$1$ góc vuông}
\loigiai{
Hình chữ nhật có 4 góc vuông.\\
\textbf{Đáp án của câu hỏi là C}
}
\end{ex}

\begin{ex} %[4H6V27]
Cho đoạn thẳng $AB$. Qua điểm $M$ nằm ngoài đoạn thẳng $AB$, ta có thể vẽ được bao nhiêu đường thẳng vuông góc với $AB$?
\choice
{$0$ đường thẳng}
{$2$ đường thẳng}
{Vô số đường thẳng}
{\True $1$ đường thẳng duy nhất}
\loigiai{
Qua một điểm cho trước, chỉ vẽ được duy nhất 1 đường thẳng vuông góc với đường thẳng đã cho.\\
\textbf{Đáp án của câu hỏi là D}
}
\end{ex}

\begin{ex} %[4H6V27]
Hình thang vuông là hình thang có
\choice
{1 góc tù}
{\True 1 cạnh bên vuông góc với hai đáy (có 2 góc vuông)}
{4 góc vuông}
{2 đường chéo vuông góc}
\loigiai{
Hình thang vuông có 1 cạnh bên vuông góc với hai đáy.\\
\textbf{Đáp án của câu hỏi là B}
}
\end{ex}

\begin{ex} %[4H6V27]
Hai kim đồng hồ lúc 3 giờ đúng tạo thành hai tia
\choice
{\True Vuông góc với nhau}
{Song song}
{Thẳng hàng}
{Trùng nhau}
\loigiai{
Lúc 3 giờ đúng hai kim vuông góc với nhau ($90^\circ$).\\
\textbf{Đáp án của câu hỏi là A}
}
\end{ex}

\begin{ex} %[4H6V27]
Trong một hình tam giác vuông cân, hai cạnh góc vuông có độ dài
\choice
{Khác nhau}
{Gấp đôi nhau}
{\True Bằng nhau}
{Không so sánh được}
\loigiai{
Tam giác vuông cân có hai cạnh góc vuông bằng nhau.\\
\textbf{Đáp án của câu hỏi là C}
}
\end{ex}

\begin{ex} %[4H6V27]
Cho hình vuông $ABCD$. Hai đường chéo $AC$ và $BD$ cắt nhau tại $O$. Tại điểm $O$ có bao nhiêu góc vuông được tạo thành bởi hai đường chéo?
\choice
{$2$ góc vuông}
{$1$ góc vuông}
{$0$ góc vuông}
{\True $4$ góc vuông}
\loigiai{
Hai đường chéo hình vuông vuông góc tại $O$ tạo thành 4 góc vuông.\\
\textbf{Đáp án của câu hỏi là D}
}
\end{ex}

\textbf{Phần 2. Trắc nghiệm đúng sai}
\setcounterpageref{ex}{0}

\begin{ex} %[4H6V27]
Xét tính đúng sai về hai đường thẳng vuông góc:
\choiceTF
{\True Hai đường thẳng vuông góc cắt nhau tạo thành các góc có số đo $90^\circ$.}
{\True Dùng ê-ke để kiểm tra và vẽ đường thẳng vuông góc.}
{Hai đường thẳng cắt nhau bất kì luôn vuông góc với nhau.}
{\True Trong hình chữ nhật, mỗi cạnh đều vuông góc với hai cạnh kề với nó.}
\loigiai{
\begin{itemize}
    \item Vuông góc tạo góc $90^\circ$: Đúng.
    \item Dùng ê-ke: Đúng.
    \item Hai đường thẳng cắt nhau có thể tạo góc nhọn/tù (không nhất thiết vuông góc). Sai.
    \item Mỗi cạnh vuông góc với 2 cạnh kề: Đúng.
\end{itemize}
\begin{itemize}
    \item a) \textbf{Đúng}.
    \item b) \textbf{Đúng}.
    \item c) \textbf{Sai}.
    \item d) \textbf{Đúng}.
\end{itemize}
}
\end{ex}

\begin{ex} %[4H6V27]
Cho hình chữ nhật $MNPQ$. Xét tính đúng sai:
\choiceTF
{\True Cạnh $MN$ vuông góc với cạnh $MQ$.}
{\True Cạnh $NP$ vuông góc với cạnh $PQ$.}
{Cạnh $MN$ vuông góc với cạnh $PQ$.}
{\True Cạnh $MQ$ vuông góc với cạnh $QP$.}
\loigiai{
\begin{itemize}
    \item $MN \perp MQ$: Đúng.
    \item $NP \perp PQ$: Đúng.
    \item $MN$ song song với $PQ$ (không vuông góc). Sai.
    \item $MQ \perp QP$: Đúng.
\end{itemize}
\begin{itemize}
    \item a) \textbf{Đúng}.
    \item b) \textbf{Đúng}.
    \item c) \textbf{Sai}.
    \item d) \textbf{Đúng}.
\end{itemize}
}
\end{ex}

\textbf{Phần 3. Trả lời ngắn}
\setcounterpageref{ex}{0}

\begin{ex} %[4H6V27]
Một hình chữ nhật có bao nhiêu góc vuông?
\shortans{4}
\loigiai{
\textbf{Đáp số:} $4$\\
Có 4 góc vuông.
}
\end{ex}

\begin{ex} %[4H6V27]
Hai đường thẳng vuông góc tạo thành bao nhiêu góc vuông tại điểm giao nhau?
\shortans{4}
\loigiai{
\textbf{Đáp số:} $4$\\
4 góc vuông.
}
\end{ex}

\begin{ex} %[4H6V27]
Một hình tam giác vuông có bao nhiêu góc vuông?
\shortans{1}
\loigiai{
\textbf{Đáp số:} $1$\\
1 góc vuông.
}
\end{ex}

\begin{ex} %[4H6V27]
Một hình chữ nhật có bao nhiêu cặp cạnh vuông góc với nhau?
\shortans{4}
\loigiai{
\textbf{Đáp số:} $4$\\
4 cặp cạnh kề nhau.
}
\end{ex}

\begin{ex} %[4H6V27]
Một hình thang vuông có mấy góc vuông?
\shortans{2}
\loigiai{
\textbf{Đáp số:} $2$\\
2 góc vuông.
}
\end{ex}

\begin{ex} %[4H6V27]
Trong một hình vuông, tính tổng số góc vuông được tạo bởi các cạnh của hình vuông và hai đường chéo cắt nhau của nó.
\shortans{8}
\loigiai{
\textbf{Đáp số:} $8$\\
4 góc vuông tại 4 đỉnh $+ 4$ góc vuông tại giao điểm hai đường chéo $= 8$ góc vuông.
}
\end{ex}
